% !TEX program = pdflatex
% Учебное пособие: как оценивать сложность рефал-функций для разных представлений объектных выражений.
% Сборка: pdflatex refal_complexity_guide.tex (из каталога nir/)
\documentclass[12pt,a4paper]{article}
\usepackage[T2A]{fontenc}
\usepackage[utf8]{inputenc}
\usepackage[russian]{babel}
\usepackage{geometry}
\geometry{margin=2.5cm}
\usepackage{amsmath}
\usepackage{booktabs}
\usepackage{listings}
\usepackage{xcolor}
\usepackage{hyperref}
\hypersetup{colorlinks=true,linkcolor=blue,urlcolor=blue}

\definecolor{codegreen}{rgb}{0,0.5,0}
\lstdefinestyle{refal}{
  language={},
  basicstyle=\ttfamily\small,
  keywordstyle=\color{blue},
  commentstyle=\color{codegreen},
  stringstyle=\color{magenta!70!black},
  breaklines=true,
  frame=single,
  numbers=left,
  numberstyle=\tiny,
  tabsize=2,
  showstringspaces=false
}
\lstset{style=refal}

\newcommand{\bigO}{\ensuremath{O}}

\title{Оценка сложности рефал-функций\\для разных представлений объектных выражений}
\author{Учебное приложение к НИР}
\date{}

\begin{document}
\maketitle

\tableofcontents
\newpage

\section{Зачем это нужно}

Интерпретатор Рефала на каждом шаге выполняет две группы действий:
\begin{enumerate}
  \item \textbf{Сопоставление} аргумента функции с образцом левой части предложения: разбиение выражения на фрагменты, связывание переменных.
  \item \textbf{Построение результата} правой части: подстановка значений переменных, вложенные вызовы, конкатенация фрагментов в одно объектное выражение.
\end{enumerate}

Стоимость этих действий в \emph{единицах работы с памятью} зависит от того, \textbf{как в памяти представлено объектное выражение}. В записке рассматривались (в упрощённом виде):
\begin{itemize}
  \item \textbf{Классическое (плоское двусвязное)} --- узел на каждый знак записи (символ или скобка), быстрые концы и конкатенация, дорогое полное копирование.
  \item \textbf{Компромиссное (подвешенные скобки)} --- узел на терм верхнего уровня, копирование верхнего уровня дешевле, но сложнее память.
  \item \textbf{Массивное} --- непрерывный блок, \(e\)-переменная как пара индексов (копирование дёшево), конкатенация часто дорогая.
  \item \textbf{Векторно-списковое} --- результат как список фрагментов вектора, конкатенация откладывается.
  \item \textbf{Rope (гипотетически)} --- дерево конкатенаций: дешёвая склейка, дороже работа с концами.
  \item \textbf{Раскрученная двусторонняя очередь} --- амортизированно дёшевые концы и конкатенация при структурном разделении.
\end{itemize}

Ниже --- \textbf{алгоритм мысленной оценки}: не формальное доказательство, а порядок рассуждений, которым можно пользоваться на экзамене или при разборе чужой функции.

\section{Кратко о синтаксисе Рефала-5 (достаточно для разбора сложности)}

\subsection{Программа, функция, предложение}

\begin{itemize}
  \item \textbf{Программа} --- набор функций.
  \item \textbf{Функция} --- имя и блок из одного или нескольких \textbf{предложений} в фигурных скобках.
  \item \textbf{Предложение}: \texttt{образец = результат;}
  \item При вызове интерпретатор просматривает предложения \textbf{сверху вниз} и берёт \textbf{первое}, образец которого сопоставим с аргументом.
\end{itemize}

Вызов функции записывают угловыми скобками: \texttt{<Имя аргумент>}.

\subsection{Объектное выражение и термы}

\textbf{Объектное выражение} --- конечная последовательность \textbf{термов} (возможно пустая).

Терм бывает:
\begin{itemize}
  \item \textbf{атомарный} (\textbf{символ} в терминологии Рефала): символ-слово (\texttt{True}), макроцифра (\texttt{42}), литера из строкового литерала (\texttt{'a'} --- один символ-литера);
  \item \textbf{скобочный}: \texttt{( ... )} --- внутри снова объектное выражение.
\end{itemize}

\subsection{Переменные образца}

\begin{itemize}
  \item \texttt{s.X} --- ровно один символ (атомарный терм).
  \item \texttt{t.X} --- ровно один терм (в том числе скобочный).
  \item \texttt{e.X} --- цепочка термов произвольной длины, в том числе пустая.
\end{itemize}

Повтор одного и того же имени переменной в образце означает \textbf{совпадение значений} (как в листинге с \texttt{Contains} ниже).

\subsection{Примеры: сопоставление (левая часть) и построение результата (правая часть)}

Ниже для \textbf{каждого} примера: какой аргумент, как интерпретатор находит первое подходящее предложение (сопоставление), из чего потом собирается результат. Порядок вычисления в правой части: сначала вычисляются вложенные вызовы \texttt{<...>} слева направо, затем конкатенация текстовых фрагментов результата в одно выражение.

\medskip
\noindent\textbf{Вызов и аргумент --- не одно и то же.}
\begin{itemize}
  \item \textbf{Вызов функции} --- целиком запись в угловых скобках \texttt{<$\ldots$>}: имя функции и переданное ей выражение.
  \item \textbf{Аргумент} --- только \textbf{объектное выражение после имени} (цепочка термов до закрывающей \texttt{>}). Имя функции в аргумент \textbf{не} входит; угловые скобки тоже не входят --- это синтаксис вызова вокруг аргумента.
\end{itemize}

\noindent Иллюстрация на одной строке (пробел после имени --- по стилю Рефала-5; его можно опускать, если однозначно):

\begin{lstlisting}
<TakeRight A : B C>
\end{lstlisting}

Здесь \textbf{вызов} --- вся строка \texttt{<TakeRight A : B C>}. \textbf{Аргумент} --- подвыражение \texttt{A : B C} (четыре терма: \texttt{A}, \texttt{:}, \texttt{B}, \texttt{C}). Левую часть предложения сопоставляют именно с \texttt{A : B C}, а не с именем \texttt{TakeRight}.

\subsubsection*{Пример 1. Пустой аргумент и константа в результате}

\begin{lstlisting}
EmptyLen {
   = 0;
}
\end{lstlisting}

\noindent\textbf{Вызов (пример):}
\begin{lstlisting}
<EmptyLen >
\end{lstlisting}
(между именем и \texttt{>} нет термов --- \textbf{аргумент пустой}.)

\textbf{Аргумент:} пустое выражение (ни одного терма).

\textbf{Сопоставление.} В образце первого предложения нет ни символов, ни переменных --- это \textbf{пустой образец}. Он сопоставим только с пустым аргументом. Переменные не связываются.

\textbf{Правая часть.} Результат --- одна макроцифра \texttt{0}. Дополнительных операций над полем зрения нет.

\subsubsection*{Пример 2. Один символ с головы и хвост}

\begin{lstlisting}
HeadEcho {
  s.H e.T = s.H (s.H) <HeadEcho e.T>;
   = ;
}
\end{lstlisting}

\noindent\textbf{Вызов (пример):}
\begin{lstlisting}
<HeadEcho A B>
\end{lstlisting}
\textbf{Аргумент первого шага:} \texttt{A B} (два символа). На рекурсии --- \texttt{<HeadEcho B>}, аргумент \texttt{B}; затем \texttt{<HeadEcho >} с пустым аргументом.

\textbf{Сопоставление второго предложения.} Слева направо: \texttt{s.H} забирает ровно один символ с \textbf{начала} --- \texttt{H$\mapsto$A}. Затем \texttt{e.T} забирает всё оставшееся до конца аргумента --- \texttt{T$\mapsto$B}. (Если бы после \texttt{s.H} был ещё фиксированный терм, граница для \texttt{e.T} сузилась бы.)

\textbf{Правая часть.} Сначала вычисляется \texttt{<HeadEcho e.T>}, т.\,е. \texttt{<HeadEcho B>}. Потом результат всего предложения --- конкатенация трёх фрагментов: \texttt{s.H}, скобочного \texttt{(s.H)} (внутри копия того же символа), и ответа рекурсии. Для \texttt{B} сработает третье предложение \texttt{ = ;} --- пустой результат. Итог вызова на \texttt{A B}: \texttt{A (A)}.

\subsubsection*{Пример 3. Фиксированные термы-якоря и одна \texttt{e}-переменная}

\begin{lstlisting}
Inside {
  'a' e.Mid 'b' = e.Mid;
}
\end{lstlisting}

\noindent\textbf{Вызов (пример):}
\begin{lstlisting}
<Inside 'a' X Y 'b'>
\end{lstlisting}
\textbf{Аргумент:} \texttt{'a' X Y 'b'} (якоря --- литеры \texttt{'a'} и \texttt{'b'}).

\textbf{Сопоставление.} Первый терм аргумента обязан совпасть с \texttt{'a'} --- да. Дальше \texttt{e.Mid} должна поглотить середину так, чтобы остаток начинался с \texttt{'b'}. При единственном вхождении \texttt{'b'} на конце получается \texttt{Mid$\mapsto$ X Y}. Если бы \texttt{'b'} встречалось дважды, без дополнительных ограничений возможны \textbf{несколько} разбиений --- перебор.

\textbf{Правая часть.} Результат --- одно значение \texttt{e.Mid}, т.\,е. копия/передача подцепочки \texttt{X Y} без якорей.

\subsubsection*{Пример 4. Повтор имени переменной в образце}

\begin{lstlisting}
PairEq {
  s.A s.B s.A = (s.A s.B);
}
\end{lstlisting}

\noindent\textbf{Вызов (пример):}
\begin{lstlisting}
<PairEq X Y X>
\end{lstlisting}
\textbf{Аргумент:} \texttt{X Y X}.

\textbf{Сопоставление.} \texttt{s.A$\mapsto$X}, \texttt{s.B$\mapsto$Y}, третий символ снова должен быть \texttt{X} --- совпадает с уже связанным \texttt{s.A}. Иначе предложение не подходит.

\textbf{Правая часть.} Скобочный терм \texttt{(s.A s.B)} --- один терм с внутренностью \texttt{X Y}.

\subsubsection*{Пример 5. Две \texttt{e}-переменные и разделитель}

\begin{lstlisting}
TakeRight {
  e.Left ':' e.Right = e.Right;
}
\end{lstlisting}

\noindent\textbf{Вызов (пример):}
\begin{lstlisting}
<TakeRight A : B C>
\end{lstlisting}
\textbf{Аргумент:} \texttt{A : B C} (не путать с \textbf{правой} частью предложения \texttt{= e.Right;} --- там \texttt{e.Right} уже \emph{имя} куска аргумента после сопоставления).

\textbf{Сопоставление.} Рефал перебирает варианты длины \texttt{e.Left}, пока не найдётся такое разбиение, что после выделенной части идёт \texttt{':'}, а за ним --- \texttt{e.Right}. Здесь единственное разбиение: \texttt{Left$\mapsto$A}, \texttt{Right$\mapsto$ B C}.

\textbf{Правая часть.} Возвращается только \texttt{e.Right} --- \texttt{B C}.

\subsubsection*{Пример 6. Таблица в скобках и рекурсия (как у \texttt{Subst})}

\begin{lstlisting}
CopyTwice {
  (e.Tab) s.X = (e.Tab) (e.Tab) s.X;
}
\end{lstlisting}

\noindent\textbf{Вызов (пример):}
\begin{lstlisting}
<CopyTwice (A B) Z>
\end{lstlisting}
\textbf{Аргумент:} \texttt{(A B) Z} --- два терма верхнего уровня: скобочный \texttt{(A B)} и символ \texttt{Z}.

\textbf{Сопоставление.} Первый терм --- скобочный; внутри образца \texttt{(e.Tab)} переменная \texttt{e.Tab} связывается с внутренностью таблицы: \texttt{Tab$\mapsto$ A B}. \texttt{s.X$\mapsto$Z}.

\textbf{Правая часть.} Конкатенация трёх кусков: две копии \texttt{(e.Tab)} и символ \texttt{s.X}. Здесь \texttt{e.Tab} в левой части \textbf{один} раз, в правой --- \textbf{два} --- типичный сигнал к вопросу о \textbf{физическом копировании} при списковом представлении (см. шаг~2 алгоритма ниже).

\section{Операции над выражением как абстрактный тип}

Чтобы оценивать сложность, достаточно отслеживать, сколько раз и в каком порядке вызываются типичные операции реализации:

\begin{enumerate}
  \item \textbf{Пусто?} --- проверка на пустое выражение.
  \item \textbf{Отделить начало} --- взять первый терм и остаток (для списков часто \(O(1)\) на указателях).
  \item \textbf{Отделить конец} --- взять последний терм и остаток (для двусвязного списка \(O(1)\), для rope обычно дороже).
  \item \textbf{Конкатенация} --- склеить два выражения в одно (у классического списка \(O(1)\), у наивного массива часто \(O(|A|+|B|)\)).
  \item \textbf{Сингуляр} --- выражение из одного терма.
  \item \textbf{Копирование значения переменной} --- когда одно и то же \(e\) (или вся таблица) используется в правой части больше раз, чем в левой: в списковых представлениях часто требуется \textbf{физическая копия} (\(O(\text{размер})\)), в массивном --- часто достаточно скопировать \textbf{диапазон} как пару индексов (\(O(1)\)).
\end{enumerate}

\section{Алгоритм оценки сложности (по шагам)}

Дана функция (или одно предложение). Нужно оценить время \textbf{одного успешного применения предложения} и/или всей функции на входе длины \(n\).

\subsection{Шаг 1. Разметить образец}

Пройти образец слева направо и записать, что происходит с аргументом:
\begin{itemize}
  \item фиксированный терм в образце --- сравнение с текущей позицией (обычно \(O(1)\) за терм, если не нужно «отматывать» длинный перебор);
  \item \texttt{s.*} --- отделить один символ с начала (или с конца, если образец устроен так --- смотрите на то, как Рефал сопоставляет; на практике важны оба конца);
  \item \texttt{t.*} --- отделить один терм (скобочный терм может потребовать прыжка по паре скобок в плоском представлении --- всё ещё часто \(O(1)\) на шаг);
  \item \texttt{e.*} --- либо однозначное «съесть всё до известной границы», либо \textbf{неоднозначность}: несколько вариантов разбиения.
\end{itemize}

\subsection{Шаг 2. Подсчитать использование переменных в правой части}

Для каждой переменной \(X\) сравнить:
\begin{itemize}
  \item \(k_L\) --- сколько раз \(X\) встречается в \textbf{левой} части;
  \item \(k_R\) --- сколько раз в \textbf{правой} части.
\end{itemize}

Если \(k_R > k_L\), то при классическом списковом представлении часто нужно \(k_R - k_L\) \textbf{копий} большого значения (типичный источник квадратики в \texttt{Subst}).

\subsection{Шаг 3. Разобрать правую часть как дерево вызовов}

Вложенные вызовы \texttt{<G ...>} сначала вычисляются, их результаты затем конкатенируются. Каждая явная конкатенация в результате --- потенциальная линейная работа для массива.

\subsection{Шаг 4. Рекурсия}

Записать рекуррентное соотношение: на сколько уменьшается размер аргумента за один вызов, сколько конкатенаций/копий на шаг.

\subsection{Шаг 5. Свести к модели представления}

Подставить стоимость операций из таблицы ниже (порядок, не константы).

\textbf{Важно.} В каждой ячейке таблицы указан порядок \textbf{одной} примитивной операции (одно отделение с конца, одна конкатенация двух кусков и т.\,д.). Оценка функции --- это, как правило, «сколько раз эта операция выполняется на всём входе» \(\times\) «что в ячейке»: например, \(\Theta(n)\) раз по \(O(1)\) из столбца «Концы» даёт \(\bigO(n)\), а не \(O(1)\).

\begin{table}[h]
\centering
\begin{tabular}{@{}lccc@{}}
\toprule
Представление & Концы & Конкатенация & Копия переменной \\
\midrule
Классическое списковое & \(O(1)\) & \(O(1)\) & \(O(N)\) \\
Компромиссное & \(O(1)\) & \(O(1)\) & \(O(N_{\text{верх}})\) \\
Массивное & \(O(1)\) & \(O(|A|+|B|)\) & \(O(1)\) \\
Векторно-списковое & \(O(1)\) & отложенная & \(O(1)\) \\
Rope & \(O(\log n)\) & \(O(1)\)/баланс & \(O(1)\) \\
\parbox[t]{3.35cm}{\raggedright Раскрученная двусторонняя очередь} & \(O(1)\) аморт. & \(O(1)\) аморт. & \(O(1)\) \\
\bottomrule
\end{tabular}
\caption{Порядки стоимости базовых операций (последняя строка --- предлагаемое в НИР представление; схематично)}
\end{table}

\subsection{Таблица и «ячейки»: что именно запоминать}

Таблица --- это не одна «клетка на всю функцию», а \textbf{справочник стоимости базовых приёмов} для каждого представления:
\begin{itemize}
  \item \textbf{Строка} --- как устроено поле зрения (классический список, массив, раскрученная двусторонняя очередь и т.\,д.).
  \item \textbf{Столбец} --- вид операции: работа с \textbf{концами} (отделить начало/конец), \textbf{конкатенация}, \textbf{копия} большого значения переменной.
\end{itemize}

\textbf{Как отвечать на вопрос про сложность.} (1)~По шагам 1--5 (подразделы «Шаг 1»--«Шаг 5» в разделе~4 выше) выписать, какие из этих операций реально возникают при \textbf{сопоставлении} и при \textbf{построении правой части} и сколько раз на один шаг / на весь вход. (2)~Для \textbf{каждого} интересующего представления подставить порядок из соответствующей \emph{ячейки} таблицы: например, «много конкатенаций подряд + массивное представление» $\Rightarrow$ столбец «Конкатенация» и строка «Массивное». (3)~Не забыть ловушки вне таблицы: неоднозначные \texttt{e.*} (перебор), рекурсия с уменьшением размера, \(k_R>k_L\) для длинных \texttt{e.*}.

Запомнить дословно все числа в таблице не обязательно; важно \textbf{смысл}: где конкатенация дешева, где дорога, где копия длинного хвата дешёва как «срез» индексов.

\section{Пять функций-примеров и разбор по шагам алгоритма}

Для каждой функции ниже кратко пройдены шаги 1--5 из раздела~4: образец \(\rightarrow\) переменные в правой части \(\rightarrow\) дерево вызовов \(\rightarrow\) рекурсия \(\rightarrow\) что взять из таблицы.

\subsection{Функция 1: длина выражения (\texttt{Length})}

\begin{lstlisting}
Length {
   = 0;
  s.First e.Rest = <Add 1 <Length e.Rest>>;
}
\end{lstlisting}

\textbf{Шаг 1 (образец).} Либо пустой аргумент (первое предложение), либо от головы: \texttt{s.First} --- один символ, \texttt{e.Rest} --- остаток. Одно отделение начала на успешное применение второго предложения.

\textbf{Шаг 2 (переменные).} \texttt{e.Rest} по одному разу в левой и правой части; лишних копий длинного хвоста из-за повторения нет.

\textbf{Шаг 3 (правая часть).} Сначала вычисляется \texttt{<Length e.Rest>}, затем \texttt{<Add 1 ...>}. Конкатенация поля зрения минимальна (результат --- число).

\textbf{Шаг 4 (рекурсия).} Длина аргумента уменьшается на~1 за вызов; всего \(\Theta(n)\) вызовов второго предложения для входа из \(n\) символов.

\textbf{Шаг 5 (таблица) --- ответ по \emph{каждой} строке.} Шагов сопоставления второго предложения \(\Theta(n)\); на шаг --- одно «отделить начало» (столбец «Концы»), копия длинного \texttt{e.Rest} не размножается (шаг~2). Слово «концы» в таблице --- про стоимость \textbf{одной} примитивной операции, итог \(=\) (число шагов) \(\times\) (ячейка «Концы»).

\begin{itemize}
  \item \textbf{Классическое списковое.} «Концы» \(O(1)\) за шаг \(\Rightarrow\) \(n\cdot O(1)=\bigO(n)\).
  \item \textbf{Компромиссное.} Для цепочки символов верхнего уровня (без дорогого углубления в скобки) то же поведение, что у списка на верхнем уровне \(\Rightarrow\) \(\bigO(n)\) при \(O(1)\) на шаг.
  \item \textbf{Массивное.} Если \texttt{e.Rest} --- срез/сдвиг индекса без физического копирования всего хвоста на каждом шаге, «Концы» \(O(1)\) \(\Rightarrow\) \(\bigO(n)\). Если же каждый шаг пересобирает массив под новый хвост, доминируют «Конкатенация»/копия \(\Rightarrow\) до \(\bigO(n^2)\) (ловушка модели, не ячейка «Концы» одна).
  \item \textbf{Векторно-списковое.} \(\Theta(n)\) шагов; при \(O(1)\) на сдвиг границы у буфера \(\Rightarrow\) \(\bigO(n)\) (отложенная конкатенация к результату-числу обычно не добавляет порядка).
  \item \textbf{Rope.} В таблице «Концы» \(O(\log n)\) \(\Rightarrow\) \(n\cdot O(\log n)=\bigO(n\log n)\) для такой рекурсии от головы.
  \item \textbf{Раскрученная двусторонняя очередь} (представление из НИР, не произвольная абстрактная двусторонняя очередь без раскрутки буферов): «Концы» \(O(1)\) амортизированно \(\Rightarrow\) \(\bigO(n)\) амортизированно по числу шагов.
\end{itemize}

\textbf{Итог.} Для пяти первых строк (без rope) при дешёвом отделении начала обычно \(\bigO(n)\); для rope --- \(\bigO(n\log n)\) по столбцу «Концы»; для массива обязательно оговорить, передаётся ли \texttt{e.Rest} срезом или копируется блок.

\subsection{Функция 2: последний символ (\texttt{Last})}

\begin{lstlisting}
Last {
  e.Init s.Last = s.Last;
}
\end{lstlisting}

\textbf{Шаг 1 (образец).} \texttt{e.Init} поглощает префикс, \texttt{s.Last} --- последний символ. Часто это заставляет механизм сопоставления «пробовать» длины \texttt{e.Init}; в худшем случае --- линейное движение по цепочке от головы (зависит от реализации Рефала).

\textbf{Шаг 2.} \texttt{s.Last} один раз слева и справа; \texttt{e.Init} в результат не входит --- копировать длинную цепочку в ответ не нужно.

\textbf{Шаг 3.} Результат --- один символ; построение тривиально.

\textbf{Шаг 4.} Рекурсии нет; один проход по предложениям.

\textbf{Шаг 5 (таблица) --- ответ по \emph{каждой} строке.} Одно применение предложения на аргументе из \(n\) символов; доминирует не рекурсия, а \textbf{стратегия сопоставления} образца \texttt{e.Init s.Last} (сколько раз дергается «Концы» при переборе границы \texttt{e.Init}).

\begin{itemize}
  \item \textbf{Классическое списковое.} При сопоставлении «от головы» (удлинение \texttt{e.Init}, проверка последнего терма) в худшем случае \(\Theta(n)\) операций с концами \(\Rightarrow\) \(\bigO(n)\). Если реализация сразу «смотрит» на хвост (двусвязность) --- \(\bigO(1)\).
  \item \textbf{Компромиссное.} Аналогично верхнему уровню: сканирование цепочки символов \(\Rightarrow\) \(\bigO(n)\) в худшем случае; доступ с конца при поддержке со стороны машины --- \(\bigO(1)\) по верхнему уровню.
  \item \textbf{Массивное.} Доступ к последнему элементу по индексу \(\bigO(1)\); если же сопоставление моделируется как многократный сдвиг начала без среза --- до \(\bigO(n)\). Конкатенация в результате не тянет порядок (один символ).
  \item \textbf{Векторно-списковое.} Обычно \(\bigO(1)\) амортизированно к последнему фрагменту + границе буфера или \(\bigO(n)\) при наивном проходе --- зависит от реализации сопоставления.
  \item \textbf{Rope.} Доступ к «последнему» символу через дерево: ячейка «Концы» \(O(\log n)\) на одно успешное сопоставление \(\Rightarrow\) \(\bigO(\log n)\) за вызов при хорошей реализации; при переборе вариантов границы --- умножить на число попыток.
  \item \textbf{Раскрученная двусторонняя очередь} (НИР). Доступ к хвосту и отделение последнего символа амортизированно \(\bigO(1)\) при сопоставлении с конца \(\Rightarrow\) \(\bigO(1)\) амортизированно за одно предложение; сопоставление только с головы --- как у списка, \(\bigO(n)\) в худшем случае.
\end{itemize}

\textbf{Итог.} Диапазон от \(\bigO(1)\) (хвост, РДО/двусвязность, явная поддержка) до \(\bigO(n)\) или \(\bigO(\log n)\) за одно предложение --- обязательно сказать, \emph{откуда} машина сопоставляет образец.

\subsection{Функция 3: неоднозначность и перебор (\texttt{Contains})}

\begin{lstlisting}
Contains {
  s.Needle e.Left s.Needle e.Right = True;
  s.Needle e.Any                   = False;
}
\end{lstlisting}

\textbf{Шаг 1 (образец).} Первое предложение: \texttt{s.Needle}, затем \texttt{e.Left}, снова тот же символ \texttt{s.Needle}, затем \texttt{e.Right}. Неоднозначность: где заканчивается \texttt{e.Left} и где второе вхождение \texttt{s.Needle} --- перебор вариантов (бэктрекинг).

\textbf{Шаг 2.} Все переменные в результате не используются (только \texttt{True}/\texttt{False}); проблема копий \texttt{e.*} в ответе не стоит.

\textbf{Шаг 3.} Правая часть --- константа; дорого только сопоставление.

\textbf{Шаг 4.} Рекурсии нет; число попыток сопоставления первого предложения может быть большим.

\textbf{Шаг 5 (таблица) --- ответ по \emph{каждой} строке.} Пусть в худшем случае \(\bigO(n^2)\) кандидатов перебора границы между \texttt{e.Left} и вторым \texttt{s.Needle}; \(T_{\text{канд}}\) --- стоимость одного кандидата по столбцу «Концы» (+ при необходимости «Конкатенация» для бэктрекинга). Итог \(\bigO(n^2\cdot T_{\text{канд}})\) до полиномов выше, если \(T_{\text{канд}}\) не константа.

\begin{itemize}
  \item \textbf{Классическое списковое.} Сдвиг границы с головы \(\Rightarrow\) \(T_{\text{канд}}=O(1)\) \(\Rightarrow\) \(\bigO(n^2)\) в типичной оценке.
  \item \textbf{Компромиссное.} Как у списка на верхнем уровне для цепочки символов \(\Rightarrow\) \(\bigO(n^2)\) при \(O(1)\) на кандидата; если перебор задевает вложенные скобки, константа в \(T_{\text{канд}}\) растёт с глубиной структуры.
  \item \textbf{Массивное.} Если каждый кандидат требует сдвига/копирования префикса массива, \(T_{\text{канд}}\) может быть \(O(n)\) \(\Rightarrow\) до \(\bigO(n^3)\); при индексах и срезах без копии --- ближе к \(\bigO(n^2)\).
  \item \textbf{Векторно-списковое.} Часто \(T_{\text{канд}}=O(1)\) амортизированно за счёт буферов \(\Rightarrow\) \(\bigO(n^2)\) по числу кандидатов.
  \item \textbf{Rope.} \(T_{\text{канд}}=O(\log n)\) по ячейке «Концы» \(\Rightarrow\) \(\bigO(n^2\log n)\) в худшем случае при том же числе кандидатов.
  \item \textbf{Раскрученная двусторонняя очередь} (НИР). \(T_{\text{канд}}=O(1)\) амортизированно при сдвиге границы у конца буфера \(\Rightarrow\) \(\bigO(n^2)\) амортизированно (как у списка, но с учётом редких дорогих ребалансировок внутри длинной серии кандидатов).
\end{itemize}

\textbf{Итог.} Таблица задаёт \(T_{\text{канд}}\); число кандидатов --- из образца и входа. Наиболее контрастно: \(\bigO(n^2)\) против \(\bigO(n^2\log n)\) (rope) и возможный \(\bigO(n^3)\) на наивном массиве.

\subsection{Функция 4: копирование таблицы (\texttt{Subst})}

\begin{lstlisting}
Subst {
  (e.Tab)            = ;
  (e.Tab) s.X e.Rest = <Lookup (e.Tab) s.X> <Subst (e.Tab) e.Rest>;
}
\end{lstlisting}

\textbf{Шаг 1 (образец).} Второе предложение: один скобочный терм с внутренностью \texttt{e.Tab}, затем \texttt{s.X}, \texttt{e.Rest}. Сопоставление с головы: отделить таблицу как один \texttt{t}-образный кусок, затем символ, затем хвост.

\textbf{Шаг 2.} \texttt{(e.Tab)} в левой части \textbf{один} раз, в правой --- в \texttt{<Lookup (e.Tab) s.X>} и снова в \texttt{<Subst (e.Tab) e.Rest>} --- \textbf{два} раза: \(k_R>k_L\) для содержимого таблицы. Это прямое попадание в столбец «Копия переменной» для списковых представлений.

\textbf{Шаг 3.} Сначала \texttt{Lookup}, затем рекурсивный \texttt{Subst}; между ними конкатенация двух результатов в поле зрения.

\textbf{Шаг 4.} За шаг укорачивается \texttt{e.Rest} на один символ; \(\Theta(n)\) шагов при длине аргумента \(n\) после таблицы.

\textbf{Шаг 5 (таблица) --- ответ по \emph{каждой} строке.} \(\Theta(n)\) рекурсивных шагов по длине \texttt{e.Rest}; на шаг --- повтор \texttt{(e.Tab)} в правой части (шаг~2) + конкатенация двух кусков + \texttt{Lookup}. Пусть \(|T|=m\) --- размер таблицы (в терминах копируемого объёма для списков).

\begin{itemize}
  \item \textbf{Классическое списковое.} «Копия переменной» \(O(m)\) на шаг \(\Rightarrow\) \(\bigO(nm)\); при \(m\sim n\) часто \(\bigO(n^2)\). «Конкатенация» мелких к \texttt{e.Rest} --- обычно не доминирует.
  \item \textbf{Компромиссное.} Копия по верхнему уровню \(O(N_{\text{верх}})\) на шаг \(\Rightarrow\) \(\bigO(n\cdot N_{\text{верх}})\); вложенность таблицы уменьшает эффективный размер копии по сравнению с классикой.
  \item \textbf{Массивное.} «Копия» таблицы как срез \(O(1)\) на шаг; \texttt{Lookup} при \(O(1)\) \(\Rightarrow\) \(\bigO(n)\). Конкатенация: если каждый шаг склеивает два массива по ячейке \(O(|A|+|B|)\), возможна \(\bigO(n^2)\) по длине накапливаемого результата --- оговорить модель.
  \item \textbf{Векторно-списковое.} Копия/срез \(O(1)\), конкатенации отложенные \(\Rightarrow\) часто \(\bigO(n)\) + одна материализация; итог зависит от политики \texttt{Lookup}.
  \item \textbf{Rope.} Копия указателей на поддеревья \(O(1)\) за передачу таблицы; конкатенация по таблице \(O(1)\); обход для \texttt{Lookup} может дать \(\bigO(n\log m)\) или \(\bigO(n)\) в зависимости от реализации --- в оценке оговорить.
  \item \textbf{Раскрученная двусторонняя очередь} (НИР). Как у классического списка по порядку копии подвыражения таблицы, если \texttt{e.Tab} материализуется при каждом вложенном вызове --- \(\bigO(nm)\) в худшем случае; если таблица передаётся структурно без полной копии буферов (модель из записки) --- уточнять по инвариантам, ближе к списку или к массиву по копии.
\end{itemize}

\textbf{Итог.} Для \texttt{Subst} решает столбец «Копия переменной» \(\times n\) шагов; массив и векторно-списковое снимают копию таблицы; rope добавляет логарифмические факторы к доступу; РДО --- как у «списковых» сценариев, если не ввести явно дешёвую передачу таблицы без полной копии.

\subsection{Функция 5: рекурсивная конкатенация результата (\texttt{Map})}

\begin{lstlisting}
Map {
   = ;
  s.X e.Rest = <F s.X> <Map e.Rest>;
}
\end{lstlisting}

\textbf{Шаг 1 (образец).} Как у \texttt{Length}: \texttt{s.X}, \texttt{e.Rest} --- одно отделение начала.

\textbf{Шаг 2.} \texttt{e.Rest} не размножается в правой части лишний раз; опасность не в копии \texttt{e.*}, а в повторяющейся \textbf{конкатенации} результатов.

\textbf{Шаг 3.} Правая часть --- конкатенация результата \texttt{<F s.X>} с результатом \texttt{<Map e.Rest>}. Порядок: сначала внутренние вызовы, потом склейка.

\textbf{Шаг 4.} \(n\) рекурсивных уровней при \(n\) символах.

\textbf{Шаг 5 (таблица) --- ответ по \emph{каждой} строке.} \(n\) уровней рекурсии; на уровень --- одно отделение начала (как у \texttt{Length}) + одна конкатенация результата \texttt{<F s.X>} с уже накопленным хвостом. Пусть \texttt{<F s.X>} за \(O(1)\) по размеру поля зрения.

\begin{itemize}
  \item \textbf{Классическое списковое.} «Конкатенация» \(O(1)\) \(\Rightarrow\) \(\bigO(n)\) суммарно; «Концы» для \texttt{s.X}/\texttt{e.Rest} тоже \(O(1)\) на шаг.
  \item \textbf{Компромиссное.} Как у списка на верхнем уровне \(\Rightarrow\) \(\bigO(n)\) при \(O(1)\) на \texttt{F} и конкатенацию верхнего уровня.
  \item \textbf{Массивное.} «Конкатенация» по ячейке \(O(|A|+|B|)\): длина хвоста растёт \(1,2,\ldots,n\) \(\Rightarrow\) \(\sum_{i=1}^n i=\bigO(n^2)\) при наивной материализации после каждого шага.
  \item \textbf{Векторно-списковое.} Конкатенация отложенная \(\Rightarrow\) \(\bigO(n)\) за накопление списка фрагментов + \(\bigO(n)\) или \(\bigO(n\log n)\) за финальную материализацию --- зависит от политики, но квадратику снимает.
  \item \textbf{Rope.} «Конкатенация» \(O(1)\) за склейку; «Концы» при необходимости смотреть в глубину дерева могут дать \(\bigO(\log n)\) на шаг \(\Rightarrow\) до \(\bigO(n\log n)\) вместе с отделением, если доминирует rope.
  \item \textbf{Раскрученная двусторонняя очередь} (НИР). «Концы» и «Конкатенация» по таблице амортизированно \(O(1)\) \(\Rightarrow\) \(\bigO(n)\) амортизированно при \(O(1)\) на \texttt{F}.
\end{itemize}

\textbf{Итог.} Доминирует столбец «Конкатенация» \(\times n\) уровней (кроме случаев отложенной склейки); массив --- главный контраст (\(\bigO(n^2)\)); rope --- возможен дополнительный \(\log n\); РДО и классический список --- линейно при дешёвой склейке.

\section{Как завтра отвечать на две присланные функции}

\begin{enumerate}
  \item Кратко описать, что делает каждое предложение (1--2 предложения на русском).
  \item Для \textbf{каждого} представления из записки: перечислить, какие из шести операций (раздел~3) возникают и сколько раз на один вызов / на весь вход; затем для каждой такой операции смотреть соответствующую \textbf{ячейку} таблицы (строка --- представление, столбец --- вид работы).
  \item Отдельно отметить «ловушки»: \(k_R > k_L\) для \texttt{e.*}; неоднозначные \texttt{e.*}; цепочка конкатенаций в правой части.
  \item Дать итог в виде \(\bigO(\cdot)\) с явными \(n\), \(m\) (размер таблицы, глубина скобок и т.д.).
\end{enumerate}

\section{Замечание о точности}

Такая оценка отражает \textbf{стоимость работы с представлением выражения}. Полное время включает также поиск предложения, встроенные функции, сборку мусора на JVM и т.д. На защите обычно достаточно честно обозначить: «оцениваю доминирующие операции над полем зрения при такой-то модели представления».

\end{document}
